% Bagean: sajarah
% Bab: biografi
% Pérangan: bertrand-russell

\documentclass[../../../include/open-logic-section]{subfiles}

\begin{document}

\olfileid{his}{bio}{rus}

\olsection{Bertrand Russell}

\olphoto{russell-bertrand}{Bertrand Russell}

Bertrand Russell dipuji minangka salah siji pangadeg
filsafat analitik modhèren. Lair tanggal 18 Mei
1872, Russell ora mung kawentar amarga karya
filsafat lan logikané, nanging uga nulis akèh
buku populèr ngenani warna-warna bidhang.
Sajeroning uripé dhèwèké uga dadi aktivis
pulitik kang semangat banget.

Russell lair ing Trellech, Monmouthshire,
Wales. Wong tuwané kalebu bangsawan Inggris.
Kekaroné wong kang mikir bébas, lan malah
kekancan karo golongan radikal ing Boston
wektu iku. Sayangé, wong tuwané Russell
séda nalika dhèwèké isih cilik, lan
Russell banjur manggon karo simbah-simbahé.
Ing kono dhèwèké diasuh kanthi ajaran
agama (bab kang banget diupayakaké wong
tuwané supaya ora kelakon). Simbah
putriné ketat banget ing kabèh prakara
moral. Nalika remaja, dhèwèké akèh-akèhé
sinau ing omah kanthi bimbingan tutor
pribadi.

Pangaruh Russell marang filsafat analitik,
mligi logika, gedhé banget. Dhèwèké
sinau matématika lan filsafat ing Trinity
College, Cambridge, lan ing kono kapengaruh
déning matématikawan lan filsuf Alfred
North Whitehead. Ing taun 1910, Russell
lan Whitehead nerbitaké jilid kapisan
\emph{Principia Mathematica}, kang mbélani
pandhangan yèn matématika bisa direduksi
dadi logika. Sawisé iku dhèwèké nerbitaké
atusan buku, ésai, lan pamflèt pulitik.
Ing taun 1950, dhèwèké menang
Bebungah Nobel Sastra.

Russell melu banget ing pulitik lan
aktivisme sosial. Nalika Perang Donya~I
dhèwèké ditangkep lan dipenjara nem
sasi amarga kagiyatan pasifis lan protès.
Nalika ana ing pakunjaran, dhèwèké
isih bisa nulis lan maca, lan ngaku
yèn pengalaman iku ``rada nyenengaké.''
Dhèwèké tetep dadi pasifis sajroning
uripé, lan dipenjara manèh amarga
melu unjuk rasa kanggo nglucuti
gaman nuklir ing taun 1961. Dhèwèké
uga slamet saka kacilakan pesawat
ing taun 1948; wong-wong kang
slamet mung sing lungguh ing
bagean kanggo ngrokok. Mula Russell
ngaku yèn uripé keslametaké
déning rokok. Russell omah-omah
kaping papat, nanging kawentar uga
amarga sesambungan katresnan ing
sanjabane palakramané. Dhèwèké
séda tanggal 2 Februari 1970
nalika umur 97 taun ing
Penrhyndeudraeth, Wales.

\begin{reading}
Russell nulis otobiografi telung bagean
kang nyakup uripé saka taun 1872--1967
\citep{Russell1967,Russell1968,Russell1969}.
Bertrand Russell Research Centre ing
McMaster University nyimpen arsip
Bertrand Russell. Delengen situs wèbé
ing \citet{Duncan2015} kanggo katrangan
ngenani jilid-jilid karya klumpukané
(kalebu indèks kang bisa ditlusuri)
lan proyèk arsip. Makalah Russell
\emph{On Denoting} \citep{Russell1905}
iku karya klasik filsafat analitik
abad kaping rong puluh.

Entri Stanford Encyclopedia of Philosophy
ngenani Russell \citep{Irvine2015}
ngemot cuplikan swara Russell kang
ngrembug hasrat lan téori pulitik.
Akèh wawancara vidéo karo Russell
kasedhiya kanthi daring. Kanggo ndeleng
dhèwèké ngrembug rokok lan kacilakan
pesawat kang dialaminé, upamané delengen
\citet{RussellND}. Sawetara karya Russell,
kalebu \emph{Introduction to Mathematical Philosophy},
kasedhiya minangka buku swara gratis
ing \citet{LibriVoxND}.
\end{reading}

\end{document}
